% Mechanical relocation generated from the preservation ledger.
% Existing prose is included byte-for-byte from preserved blocks.

\begingroup
\hbadness=10000
\subsection{Numerical optimization and local uncertainty}

The current point-estimation implementation minimizes $Q_w$ with gradient-based
SciPy optimizers such as BFGS or L-BFGS-B. Gradients are supplied by JAX automatic
differentiation. Because the objective need not be globally convex, convergence
to the global optimum is not guaranteed; initialization, stopping tolerances, and
local curvature should therefore be inspected.

When requested, the software computes the Hessian of $Q_w$ at the optimum and
uses its inverse as a local covariance approximation. For ML this is an asymptotic
likelihood-based approximation; for MAP it is a Laplace-style local approximation
to posterior curvature. These standard errors are not equivalent to a complete
posterior distribution and may be unreliable for weakly identified, asymmetric,
bounded, or multimodal problems.



\subsection{Numerical Optimization, Scaling, and Precision}
\label{sec:reduced-od-numerical-contract}

The current gravity estimator evaluates a JAX value-and-gradient kernel and
uses SciPy L--BFGS--B by default.  The execution setting
\texttt{optimizer\_maxls} controls the maximum number of line-search steps
(20 by default).  An optional Biogeme TR--BFGS pilot is selected explicitly
through the gravity optimizer-comparison interface; it is not a default
dependency or a replacement for the SciPy estimator.  Optimizer and progress
settings are execution metadata and are excluded from the scientific model
fingerprint.

The raw parameter vector is mapped to physical demand and likelihood
parameters by the declared \texttt{GravityModelSpecification}.  Positive
quantities and other domain restrictions are represented by those transforms.
The default SciPy gravity fit is not a bounded optimization in raw
coordinates.  The isolated Biogeme pilot can receive explicit bounds, but
those bounds must be treated as part of that pilot's declared configuration,
not assumed to apply to the default estimator.

Convergence certification uses a scale-aware Dennis--Schnabel gradient rather
than SciPy's success flag alone.  For raw parameter $\eta_i$, objective $Q$,
typical parameter scale $\operatorname{typx}_i>0$, and objective floor
$\operatorname{typf}>0$,
\[
 g^{\mathrm{scaled}}_i
 =\frac{|\nabla_i Q|\max(|\eta_i|,\operatorname{typx}_i)}
        {\max(|Q|,\operatorname{typf})},
 \qquad
 \|g^{\mathrm{scaled}}\|_\infty
 =\max_i g^{\mathrm{scaled}}_i .
\]
The configuration records \texttt{scaled\_gradient\_tolerance},
\texttt{typical\_objective\_scale} (\texttt{typf}), and either a scalar or
complete vector of \texttt{typical\_parameter\_scales} (\texttt{typx}).
Scales are finite and strictly positive; case studies should choose them from
the initial objective and the natural units or prior scales of each parameter,
not from the final fitted values.  The result and manifest retain the resolved
scales and their provenance.

A fit is accepted as numerically converged only when the optimizer reports
success and the scaled-gradient infinity norm is no larger than the declared
tolerance.  If L--BFGS--B stops on relative objective reduction while the
scaled gradient remains large, the termination message is preserved but the
fit is reported as not converged.  An iteration or time-budget stop is
resumable when its identity-compatible checkpoint is complete, but is not an
accepted estimate.

The final evaluation also records the objective and gradient dtypes,
$\operatorname{spacing}(Q)$ near the objective, the reduction in objective
between the last two accepted iterates, and whether the requested objective
tolerance is below the available floating-point resolution.  These are
diagnostics: the package does not silently promote float32 to float64 and does
not reject a fit solely because the tolerance is finer than the spacing.  The
arithmetic follows the dtype of the supplied feature and operator arrays.

Each fit report should therefore be read as three separate statements:
\begin{description}
  \item[Optimizer termination] SciPy or the selected pilot's status, success
        flag, options, and message;
  \item[Numerical convergence] finite objective and gradients together with
        the scaled-gradient criterion and precision diagnostics;
  \item[Scientific admissibility] a case-owner decision based on adequacy,
        identification, prior sensitivity, and validation.
\end{description}
\endgroup
